\chapter{Introducción}\label{chap:introduccion}

Sostener un objeto sin tocarlo, solo con fuerzas magnéticas, es la idea que comparten todos los sistemas de levitación magnética. Es también un problema de control exigente. Un imán suspendido bajo una bobina permanece en su lugar mientras la fuerza magnética compense exactamente su peso, y una perturbación mínima basta para romper esa compensación. Esta tesis estudia cómo sostener y desplazar un imán en esas condiciones cuando, además, el imán roza con la guía que lo mantiene alineado.

El problema reúne tres dificultades que rara vez aparecen juntas en un mismo dispositivo de laboratorio. La primera es un equilibrio inestable. La segunda es una fricción que alterna entre retener al imán y soltarlo de golpe. La tercera es un actuador que no responde a los cambios pequeños del voltaje que se le ordena. Este capítulo presenta el contexto del que surge el problema, la motivación para abordarlo con un controlador lineal de estructura clásica, la literatura en la que se apoya y la formulación precisa que se resuelve en el resto del trabajo.

\section{Contexto del problema}

Los sistemas de levitación magnética tienen aplicaciones prácticas variadas, y su complejidad ha llevado a investigadores e ingenieros a explorar posibilidades aún no aprovechadas. El ejemplo más visible son los trenes de levitación magnética de alta velocidad~\cite{mingda}, el medio de transporte terrestre más relevante que emplea esta tecnología, junto con otros sistemas de transporte, muchos de ellos todavía en fase experimental~\cite{hyung}. La levitación se emplea también en la suspensión magnética activa, donde un control no lineal reduce la vibración del rotor y el consumo de energía~\cite{charara}; en los cojinetes magnéticos, que prolongan la vida útil de las máquinas al reducir el desgaste~\cite{debarghya}; en dispositivos médicos de alta precisión, como las bombas de asistencia ventricular~\cite{gang1}; y en microrrobótica, donde se ha extendido a la teleoperación en ambientes con restricciones complejas~\cite{mir}.

El costo inicial de la levitación magnética es mayor que el de una solución puramente mecánica. Sus beneficios a largo plazo ---menor mantenimiento, ahorro de energía y mayor seguridad--- la hacen, sin embargo, cada vez más competitiva en industrias como la aeroespacial~\cite{bryton}, la médica~\cite{song} y la biomédica~\cite{sajjad}. Su implementación requiere electroimanes de alta calidad, sensores precisos y conocimiento especializado para su diseño y mantenimiento. La miniaturización de los sensores, los avances de la electrónica de potencia, la mejora de los materiales magnéticos y la producción a gran escala han reducido esos costos. La demanda de menor consumo de energía, menor mantenimiento y menor huella ecológica, impulsada por programas internacionales como el Programa de las Naciones Unidas para el Medio Ambiente, ha fomentado además la inversión en esta área, y es previsible que la tecnología llegue a aplicaciones cotidianas como bombas hidráulicas, vehículos eléctricos y electrodomésticos.

Detrás de esta variedad de aplicaciones hay un mismo principio físico. La levitación magnética se basa en las fuerzas de atracción o de repulsión entre los campos de imanes permanentes y electroimanes, y la manera en que se dispone el electroimán determina el carácter del problema de control. El sistema de levitación magnética (SLM) ECP-730 de \emph{Educational Control Products}~\cite{730}, el dispositivo de laboratorio que se estudia en esta tesis, permite comparar las dos disposiciones con el mismo equipo. Un imán permanente en forma de disco se desliza a lo largo de una varilla de vidrio vertical, entre dos bobinas. En la \emph{configuración repulsiva}, la bobina inferior empuja al imán hacia arriba; en la \emph{configuración atractiva}, la bobina superior lo atrae.

Las dos configuraciones se comportan de manera opuesta ante una perturbación. En la repulsiva, si el imán sube un poco, se aleja de la bobina que lo empuja, la fuerza disminuye y el peso lo devuelve a su lugar; si baja, la fuerza aumenta y lo empuja de vuelta. El equilibrio se restablece por sí solo, y el sistema es estable en lazo abierto. En la atractiva ocurre lo contrario. Si el imán sube un poco, se acerca a la bobina que lo atrae, la fuerza crece más que el peso y el imán sigue subiendo hasta chocar con ella; si baja, la fuerza disminuye y el imán cae. Cualquier desviación, por pequeña que sea, se amplifica. Sin un controlador que corrija continuamente el voltaje de la bobina, el imán no puede permanecer en ninguna posición intermedia. En esta configuración todos los puntos de equilibrio son inestables, y estabilizarlos es la primera tarea del controlador.

Las aplicaciones con electroimanes requieren, por tanto, técnicas de control para mantener la separación entre los objetos. Una vez garantizada la estabilidad, el objetivo pasa a ser controlar el movimiento con mayor precisión~\cite{hyung}. Para evitar los desplazamientos laterales del objeto levitante, algunos dispositivos ---el ECP-730 entre ellos--- introducen guías mecánicas. Esa solución reintroduce el contacto físico que la levitación pretendía eliminar y, con él, la fricción.

La fricción entre el imán y la varilla tiene una consecuencia que conviene describir con cuidado, porque gobierna buena parte de los resultados de esta tesis. Mientras el imán está quieto, la fricción estática lo retiene aunque la fuerza neta sobre él no sea nula, siempre que esa fuerza no supere un umbral. Cuando lo supera, el imán se despega y la fricción cae a un valor menor, de modo que el movimiento arranca de forma brusca. Es la misma experiencia de empujar un mueble pesado: hace falta más fuerza para ponerlo en marcha que para mantenerlo en movimiento. La alternancia entre intervalos de reposo y deslizamientos repentinos se conoce como \emph{atascamiento-deslizamiento} (\emph{stick-slip}). En un sistema de posicionamiento, el atascamiento hace que el imán quede detenido cuando cambia la posición de referencia, y para despegarlo el controlador debe suministrar una señal suficiente para vencer la fricción estática y restablecer el movimiento.

A la fricción se suma una no linealidad del actuador. El voltaje que llega a la bobina no sigue fielmente al voltaje que ordena el controlador. Existe una banda de valores, llamada \emph{zona muerta}, dentro de la cual los cambios de la orden no producen ningún cambio en la salida. En la mayoría de los sistemas electromecánicos esa banda está centrada en el voltaje cero. En un levitador, cada posición requiere su propio voltaje de equilibrio para sostener al imán, y la zona muerta se desplaza con él. Su centro depende, entonces, de la posición en que opera el sistema.

El problema que aborda esta tesis queda así delimitado. Se trata de controlar la posición de un imán en la configuración atractiva del ECP-730, inestable en lazo abierto, en presencia de la no linealidad de la fuerza electromagnética, de la fricción con atascamiento-deslizamiento y de una zona muerta cuyo centro se mueve con el punto de operación.

\section{Motivación}

Un controlador lineal de estructura sencilla, diseñado con las herramientas del control clásico, tiene ventajas que pesan en la práctica. Su análisis se apoya en los diagramas de Bode y de Nyquist y en el lugar de las raíces, sus márgenes de estabilidad tienen una interpretación directa y su implementación en un procesador digital es inmediata. La pregunta es si una estructura así basta para un sistema que combina inestabilidad y no linealidades no suaves, o si el problema exige estrategias más elaboradas.

Dos trabajos previos, realizados en dispositivos estables, sugieren que la respuesta depende de cómo se trate la acción integral. Hernández Alcántara~\cite{diana_tesis} estudió la configuración repulsiva del ECP-730, identificó los parámetros del modelo que emplea esta tesis y atribuyó las diferencias entre el modelo y el dispositivo real a la fricción y a la zona muerta. De ahí concluyó que, para reducir sus efectos, el sistema de control debe proporcionar alta ganancia en estado estacionario. En un trabajo posterior, Hernández Alcántara et al.~\cite{alcantara} diseñaron un lazo externo con acción integral. Con él obtuvieron mejores respuestas transitorias y menores errores en estado estacionario, pero la interacción entre la acción integral y la fricción provocó oscilaciones de tipo ciclo límite, que eliminaron con un esquema de control conmutado.

El segundo trabajo apunta en una dirección complementaria. Pérez-Gómez~\cite{perez} estudió el control de posición de un motor de corriente directa (CD) con zona muerta y fricción, y comparó tres estrategias. La cancelación de la zona muerta mediante su inversa produjo oscilaciones de alta frecuencia (\emph{chattering}), y el control conmutado ofreció un error muy bajo a costa de un análisis, un diseño y una implementación complejos. El controlador lineal con doble acción integral y compensador de adelanto (PII), en cambio, se implementó en tiempo real con un error en estado estacionario y un sobreimpulso bajos, y conservó la sencillez del diseño clásico.

La doble acción integral parece, entonces, reunir lo que piden ambos antecedentes. Un segundo integrador eleva la ganancia a baja frecuencia, que es la alta ganancia en estado estacionario que recomendaba Hernández Alcántara. Además, mientras el imán permanece atascado, la acción integral acumula el error y modifica el voltaje hasta despegarlo, lo que permite esperar una salida más rápida de la zona muerta y de la fricción estática. Ninguno de los dos trabajos, sin embargo, probó esa idea en un sistema inestable. En la configuración atractiva, el controlador no solo debe vencer la fricción y la zona muerta, sino estabilizar al mismo tiempo un equilibrio que se pierde en fracciones de segundo. Extender el PII a esta configuración, y medir qué aporta realmente su segunda integración, es la motivación central de este trabajo.

A esta motivación se suma una restricción práctica. En este trabajo el voltaje del actuador se limita al intervalo de 0 a 3{,}5~V, y el voltaje que sostiene al imán crece rápidamente al alejarlo de la bobina. Un controlador útil debe operar dentro de ese intervalo, lo que obliga a preguntarse en qué posiciones puede el sistema moverse de verdad, y no solo sostenerse.

\section{Revisión bibliográfica}

La literatura relevante para este problema se organiza en tres frentes: el control de los levitadores magnéticos, el modelado y la simulación de la fricción, y el tratamiento de la zona muerta en lazos de control. Los antecedentes directos de esta tesis, ya presentados en la sección anterior, se describen con más detalle en el capítulo~\ref{Chap:Antecedentes}.

\subsection{Control de sistemas de levitación magnética}

El levitador magnético es un banco de pruebas clásico para el control no lineal, y la variedad de estrategias propuestas lo refleja. Al-Muthair y Zribi~\cite{al} aplicaron control por modos deslizantes. Charara et al.~\cite{charara} diseñaron un control no lineal para un levitador sin premagnetización. Zhang y Suyama~\cite{feifei} emplearon la linealización exacta por retroalimentación, y El Hajjaji y Ouladsine~\cite{el} combinaron el modelado del sistema con un control no lineal. Zhao et al.~\cite{zhao} desarrollaron una herramienta de control no lineal en el espacio de fases, y Khimani et al.~\cite{khimani} implementaron un control no lineal por retroalimentación de alto desempeño. Naz et al.~\cite{neelma} propusieron un control por modos deslizantes con datos muestreados y un estimador basado en el filtro de Kalman extendido, y Malik y Obaid~\cite{maaz} diseñaron un control no lineal con datos muestreados para el propio ECP-730. En el terreno del control robusto, Patil et al.~\cite{mukesh} implementaron en tiempo real un controlador basado en la teoría de retroalimentación cuantitativa (QFT, \emph{Quantitative Feedback Theory}) con prefiltro, y Su y Li~\cite{kuo1} propusieron un control supervisor basado en modelos difusos.

Dos observaciones atraviesan estos trabajos. La primera es que no existe un modelo único de la fuerza electromagnética. Cada autor la expresa de una forma distinta, condicionada por la configuración del sistema y por la manera, analítica o experimental, en que se obtuvo; el capítulo~\ref{chap:modelo} compara esas formas y justifica la que se adopta aquí. La segunda es que varios de ellos~\cite{khimani,mukesh,neelma} no incluyen la zona muerta en su modelado, ya sea porque la ignoran o porque la tratan como dinámica no modelada o como incertidumbre no estructurada.

\subsection{Fricción y atascamiento-deslizamiento}

La fricción entre superficies en contacto se describe habitualmente como la suma de varias componentes. La fricción estática impide el movimiento mientras la fuerza aplicada no supere un valor máximo; la fricción de Coulomb tiene magnitud constante y se opone a la velocidad; la fricción viscosa es proporcional a la velocidad; y el efecto Stribeck describe la disminución de la fricción al pasar del reposo a velocidades pequeñas~\cite{olsson,armstrong}. Armstrong-Hélouvry~\cite{armstrong} analizó el atascamiento-deslizamiento en el movimiento a baja velocidad y sus consecuencias para el control, y Berman et al.~\cite{berman} estudiaron el origen físico de los distintos mecanismos que lo producen.

Simular estos modelos presenta una dificultad propia, porque la fuerza de fricción es discontinua en la velocidad cero~\cite{haessig}. El método de Karnopp~\cite{karnopp} la resuelve definiendo una banda pequeña de velocidades que se tratan como reposo. Dentro de esa banda, el cuerpo permanece detenido mientras la fuerza neta no supere la fricción estática máxima, y se despega en cuanto la supera. Su sencillez explica su uso en aplicaciones prácticas. Se ha empleado para identificar la fricción en actuadores electroneumáticos~\cite{ravanbod,ravanbod1}, en péndulos~\cite{bicakci} y en válvulas de control~\cite{romano}, y Liu et al.~\cite{liu} propusieron un método estadístico para determinar la banda de velocidad nula. Esta tesis lo adopta para simular el contacto entre el imán y la varilla.

La literatura de fricción documenta también el fenómeno que resultará central en este trabajo. En los sistemas de posicionamiento, la combinación de fricción estática y acción integral produce un ciclo límite conocido como \emph{hunting}~\cite{olsson,armstrong}. Mientras el cuerpo está detenido con un error distinto de cero, el integrador acumula ese error y modifica la señal de control hasta forzar un deslizamiento; el cuerpo cruza la referencia, se detiene del otro lado y el proceso se repite. Hernández Alcántara et al.~\cite{alcantara} observaron este mecanismo en la configuración repulsiva del ECP-730.

\subsection{Zona muerta en lazos de control}

La zona muerta es una de las no linealidades no suaves más frecuentes en los actuadores, y su compensación tiene una tradición propia. Tao y Kokotović~\cite{gang} establecieron el control adaptable de plantas con zonas muertas desconocidas. A partir de ahí se han propuesto representaciones cada vez más generales: zonas muertas no simétricas~\cite{salim}, entradas con zona muerta y dirección de control desconocida~\cite{huanqing}, zonas muertas difusas~\cite{zhi}, zonas muertas inciertas en motores de corriente directa~\cite{nizar} y esquemas de compensación inversa difusa y adaptable~\cite{guanyu}. Estos modelos y los resultados experimentales que los acompañan representan solo en parte el comportamiento real de los actuadores.

En un motor de posición angular, la zona muerta se ubica alrededor de un punto fijo, porque el único voltaje de equilibrio es cero. Pérez-Gómez~\cite{perez} la modeló como una no linealidad \emph{dura} ubicada entre los subsistemas eléctrico y mecánico. En un levitador, en cambio, la zona muerta se centra en el voltaje de equilibrio de cada posición y se desplaza con él, un caso que la literatura revisada trata poco.

\subsection{Lo que falta por resolver}

La revisión deja un hueco preciso. Los controladores de levitadores reportados se diseñan, en su mayoría, sin fricción con la guía ni zona muerta. Los trabajos que sí las consideran se realizaron en dispositivos estables en lazo abierto, ya sea el ECP-730 en configuración repulsiva~\cite{diana_tesis,alcantara} o un motor de corriente directa~\cite{perez}. Queda abierto si un controlador lineal con doble acción integral estabiliza la configuración atractiva en presencia de ambas no linealidades y, sobre todo, si su segunda integración reduce el atascamiento y el ciclo límite que la acción integral provoca, como sugieren los antecedentes. Esta tesis responde esas dos preguntas.

\section{Formulación del problema}

El movimiento vertical del imán en la configuración atractiva obedece la segunda ley de Newton,
\begin{equation}
m\ddot{x} = \frac{u}{b\,(a-x)^4} - f_{\text{fr}} - mg .
\label{eq:intro_modelo}
\end{equation}
La posición $x$ del imán se mide en centímetros sobre la varilla y crece hacia la bobina superior, de modo que $a-x$ es la distancia entre ambos. El primer término de la ecuación de movimiento~\eqref{eq:intro_modelo} es la fuerza electromagnética que da el fabricante~\cite{parks}. Esa fuerza es proporcional al voltaje $u$ aplicado a la bobina y crece con la cuarta potencia inversa de la distancia, y esta dependencia es la que vuelve inestables los equilibrios. Los parámetros $a$ y $b$ y la masa $m$ del imán fueron identificados experimentalmente por Hernández Alcántara~\cite{diana_tesis}. El término $f_{\text{fr}}$ reúne la fricción con la varilla: es viscosa mientras el imán desliza y estática, hasta un valor máximo $F_s$, mientras está detenido. El último término es el peso del imán.

El voltaje $u$ que llega a la bobina pasa antes por dos restricciones del actuador. Está acotado al intervalo de 0 a 3{,}5~V y atraviesa una zona muerta centrada en el voltaje de equilibrio de la posición actual. El controlador solo dispone de la medición de la posición.

Con este modelo, el problema consiste en diseñar, sobre el modelo linealizado alrededor de un punto de equilibrio, un controlador lineal $C(s)$ que lleve al imán de una posición a otra y siga referencias variables en el tiempo dentro del intervalo que permite el actuador, sin compensar explícitamente la fricción ni la zona muerta. El controlador que se propone es un PII, y su desempeño se evalúa por simulación sobre el modelo no lineal completo.

La hipótesis que guía el trabajo tiene dos partes. La primera es que un PII diseñado sobre el modelo linealizado estabiliza la posición del levitante en presencia de la fricción estática y de la zona muerta del actuador, sin necesidad de compensarlas. La segunda es que, por la alta ganancia que proporciona a baja frecuencia, reduce el tiempo de atascamiento y la amplitud de las oscilaciones respecto a un controlador proporcional-integral (PI), con una sola acción integral. El capítulo~\ref{chap:planteamiento} enuncia formalmente el problema y la hipótesis, y el capítulo~\ref{chap:objetivos} los traduce en objetivos.

\section{Contribuciones}

La primera contribución es un controlador PII para la configuración atractiva del ECP-730, inestable en lazo abierto. Las simulaciones muestran que estabiliza el sistema en presencia de fricción estática y zona muerta, en regulación y en seguimiento de referencias senoidales y trapezoidales, sin compensar explícitamente ninguna de las dos no linealidades.

La segunda es la determinación del intervalo en que el sistema puede operar con el actuador disponible. Con un voltaje máximo de 3{,}5~V, el imán puede sostenerse en $-4$~cm, pero no desplazarse desde ahí; los cambios de referencia son viables a partir de unos $-3{,}6$~cm.

La tercera es una evaluación directa de la segunda acción integral. La comparación con un controlador PI y un barrido de 256 variantes de ambos controladores muestran que la segunda integración no reduce la amplitud del ciclo límite ni la duración de los atascamientos. La primera parte de la hipótesis se confirma y la segunda se rechaza.

La cuarta es la explicación del mecanismo del ciclo límite. La fricción estática equivale a una banda muerta de voltaje alrededor del equilibrio, y la acción integral obliga al controlador a romperla una y otra vez. La amplitud del ciclo es aproximadamente inversamente proporcional a la ganancia del controlador, y la duración de cada atascamiento la fija el término integral simple.

\section{Organización de la tesis}

El capítulo~\ref{Chap:Antecedentes} resume los trabajos previos en los que se basa esta tesis. El capítulo~\ref{chap:planteamiento} plantea el problema, la hipótesis y la justificación, y los capítulos~\ref{chap:objetivos} y~\ref{chap:metodologia} presentan los objetivos y la metodología. El capítulo~\ref{chap:marco} reúne los conceptos teóricos que se emplean en el resto del trabajo. El capítulo~\ref{chap:modelo} desarrolla el modelo del sistema, su linealización y el diseño del controlador PII. El capítulo~\ref{chap:resultados} presenta los resultados de simulación: el modelo no lineal sin y con atascamiento-deslizamiento, la regulación, el seguimiento, el análisis del ciclo límite y la configuración estable. El capítulo~\ref{chap:conclusiones} expone las conclusiones, la evaluación de la hipótesis y el trabajo futuro. Los anexos describen la implementación de las simulaciones y su verificación.
